\documentclass{notaaula}

        \titulonota{ENEM: revisão integrada}
        \creditos{Turma 3A · Ciências da Natureza para o ENEM · Novembro/2026 · Prof. Josué Cavalcante}

        \begin{document}
        \cabecalho

        \begin{sobre}
        Revisão integral por habilidades: interpretação de gráficos, unidades, modelos e situações-problema de Física, Química e Biologia.
        \end{sobre}

        \section{Método de resolução}

\begin{copiar}{Ciclo em cinco passos}
\begin{itens}
\item Leia o comando antes do texto longo e identifique a ação pedida.
\item Marque grandezas, unidades, tendência do gráfico e condição do sistema.
\item Traduza o contexto para um princípio: conservação, proporcionalidade, equilíbrio ou mecanismo biológico/químico.
\item Elimine alternativas incompatíveis antes de calcular com precisão.
\item Faça verificação dimensional e de ordem de grandeza.
\end{itens}
Se uma questão travar, registre a dúvida, avance e retorne. Tempo também é recurso de prova.
\end{copiar}

\section{Folha de relações essenciais}

\begin{copiar}{Física}
\[
  v=\frac{\Delta s}{\Delta t},\quad
  F=ma,\quad
  E_c=\frac{mv^2}{2},\quad
  E_{pg}=mgh,
\]
\[
  P=\frac{E}{\Delta t}=Ui,\quad
  U=Ri,\quad
  v_{\mathrm{onda}}=\lambda f,
\]
\[
  Q=mc\Delta T,\quad
  p=\frac{F}{A},\quad
  \rho=\frac{m}{V}.
\]
Fórmula sem modelo não basta: declare sistema, sentido positivo e unidade.
\end{copiar}

\begin{copiar}{Química e Biologia}
Em Química, conserve átomos e carga; use $n=m/M$, concentração e proporção estequiométrica.
Em Biologia, acompanhe fluxo de matéria, energia e informação: metabolismo, hereditariedade,
regulação e relações ecológicas. Em gráficos, diferencie valor absoluto de taxa e correlação de causalidade.
\end{copiar}

\section{Erros que custam pontos}

\begin{dica}
\begin{itens}
\item esquecer prefixos: m = $10^{-3}$, k = $10^3$, M = $10^6$;
\item usar minutos com potência em kW sem converter para horas;
\item confundir calor com temperatura ou potência com energia;
\item ignorar grupo de controle, dose ou escala dos eixos;
\item escolher alternativa verdadeira, mas que não responde ao comando.
\end{itens}
\end{dica}

\section{Mini-simulado misto}

\begin{exercicios}[Questões de revisão]
\begin{questoes}
\item Um carro aumenta a rapidez de $10$ para $20\un{m/s}$. Sua energia cinética
  \begin{alternativas}\item dobra\item triplica\item quadruplica\item não muda\end{alternativas}
\item Um aparelho de $500\un{W}$ funciona $2\un{h}$. O consumo é
  \begin{alternativas}\item $\dec{0,25}\un{kWh}$\item $1\un{kWh}$\item $2\un{kWh}$\item $1000\un{kWh}$\end{alternativas}
\item Uma onda de $v=300\un{m/s}$ e $f=150\un{Hz}$ tem $\lambda$
  \begin{alternativas}\item $\dec{0,5}\un{m}$\item $2\un{m}$\item $45\un{m}$\item $450\un{m}$\end{alternativas}
\item Na neutralização entre ácido e base, uma evidência relevante é
  \begin{alternativas}\item conservação de átomos e formação de produtos\item desaparecimento de massa\item criação de carga\item ausência de interação\end{alternativas}
\item Em uma cadeia alimentar, a biomassa disponível tende a diminuir nos níveis superiores devido
  \begin{alternativas}\item à criação de energia\item às perdas energéticas metabólicas\item à falta de matéria nos produtores\item ao fim da respiração\end{alternativas}
\item Uma amostra passa por três meias-vidas. A fração restante é
  \begin{alternativas}\item $1/2$\item $1/3$\item $1/6$\item $1/8$\end{alternativas}
\item Um gráfico cresce de 20 para 30 unidades. O aumento percentual sobre o valor inicial é
  \begin{alternativas}\item $10\%$\item $20\%$\item $50\%$\item $150\%$\end{alternativas}
\item Uma medida reduz a corrente na transmissão sem mudar a potência entregue. Ela deve
  \begin{alternativas}\item reduzir a tensão\item elevar a tensão\item elevar a resistência dos cabos\item zerar a energia\end{alternativas}
\item Em experimento sobre fertilizante, o melhor grupo de controle é formado por plantas
  \begin{alternativas}\item de outra espécie e ambiente\item sem fertilizante, nas mesmas condições\item com dose desconhecida\item sem medição\end{alternativas}
\item A alternativa mais científica diante de um resultado inesperado é
  \begin{alternativas}\item ocultá-lo\item alterar os dados\item revisar método e repetir medidas\item declarar lei universal\end{alternativas}
\item Explique como a análise dimensional pode eliminar uma alternativa.
\item Escolha uma questão errada do seu portfólio e registre: conceito, erro, correção e pista para reconhecer padrão semelhante.
\end{questoes}
\gabarito{1-c; 2-b; 3-b; 4-a; 5-b; 6-d; 7-c; 8-b; 9-b; 10-c.}
\end{exercicios}


\end{document}
